\chapter{Logic Tensor Networks}
\label{ch:ltn}

\Cref{ch:attacks-ddos} ended with candidate axioms written as implications
over predicates. The predicates have no thresholds, so as they stand they are
neither true nor false. \Cref{ch:learning-paradigms} placed the thesis in the
family that puts logic in the loss. This chapter presents the framework of
that family that gives the axioms a meaning and lets a network learn under
them. The framework is Logic Tensor
Networks \citep{badreddine2022ltn}. The chapter describes its language and its
semantics, the way learning is posed as an optimisation problem, what the
framework can and cannot do, and why the thesis chose it. It ends by writing
that knowledge in the language of the framework.

% ================================================================
% 5.1 Real Logic: syntax
% ================================================================
\section{Real Logic: syntax}
\label{sec:real-logic-syntax}

Real Logic is the many-valued first-order language of Logic Tensor Networks.
Its syntax is the usual one for first-order logic, with domains added as
types \citep{badreddine2022ltn}. This section covers the symbols and the
formulas. Their meaning comes from the grounding of the next section.

\subsection{Constants, variables, functions, predicates}
\label{sec:rl-symbols}

The language has five kinds of symbol: constants, variables, function symbols,
predicate symbols and domain symbols \citep{badreddine2022ltn}. Domains are
the types of the language. Every variable ranges over a domain. Every function
and every predicate has input domains, and a function has an output domain as
well. Terms are built from constants, variables and function symbols in the
usual way. An atomic formula is an equality between terms or a predicate
applied to terms of the right domains.

The axioms of \cref{ch:attacks-ddos} speak of three objects, a flow $f$, a
window of flows $w$ and a source $s$. The knowledge base of the thesis has a
single domain instead, the units that the classifier labels. The classes and
the conditions of the axioms are predicates over that domain.
\Cref{tab:rl-symbols} sets these objects against the symbols of the
language.

\begin{table}[htbp]
\centering
\footnotesize
\setlength{\tabcolsep}{3pt}
\caption{Objects of \cref{ch:attacks-ddos} as symbols of Real Logic. The
mapping is ours.}
\label{tab:rl-symbols}
\begin{tabular}{@{}p{2.6cm}p{4.9cm}p{4.8cm}@{}}
\toprule
Object & Symbol & Grounding \\
\midrule
Flow, window, source
  & One domain of classified units, with variable $x$, see \cref{sec:predicate-design}
  & Sets of real feature vectors \\
Classes
  & Predicates $\mathrm{Normal}$, $\mathrm{DoS}$, $\mathrm{Probe}$, $\mathrm{R2L}$, $\mathrm{U2R}$
  & Learned, a network with output in $[0,1]$ \\
Conditions of the axioms
  & Predicates such as $\mathrm{HighRate}$ or $\mathrm{ManySrc}$
  & Computed from the features, with output in $[0,1]$ \\
\bottomrule
\end{tabular}
\end{table}

We expect to need no function symbols. The features of a flow, a window or a
source are computed before the logic is applied, so the individuals are
already vectors.

Typing needs care. In \cref{tab:candidate-axioms} the class $\mathrm{DoS}$
takes a window in most axioms and a flow in the slow DoS axiom, while
$\mathrm{Probe}$ and $\mathrm{R2L}$ take a source. A predicate of Real Logic
has fixed input domains. The thesis moves all of them to one domain, the unit
that the classifier labels, and \cref{sec:predicate-design} gives the reason.

\subsection{Connectives and quantifiers}
\label{sec:rl-connectives}

Formulas are built with negation, conjunction, disjunction, implication and
biconditional, and with the universal and the existential quantifier
\citep{badreddine2022ltn}. The chapter writes them as $\neg$, $\wedge$,
$\vee$, $\rightarrow$, $\leftrightarrow$, $\forall$ and $\exists$. A
quantifier binds a variable, and a formula that still has free variables is
open.

The axioms of \cref{tab:candidate-axioms} leave the quantifier implicit. Read
as formulas of Real Logic they are universal. Over the domain of units, the
SYN flood axiom becomes
\[
\forall x\,\bigl(\mathrm{SynExcess}(x)\wedge\mathrm{Incomplete}(x)\rightarrow
\mathrm{DoS}(x)\bigr),
\]
where the two conditions read the statistics of the window of $x$.
At this point the formula says nothing about truth. That is the work of the
grounding.

% ================================================================
% 5.2 Grounding
% ================================================================
\section{Grounding (semantics in the tensor domain)}
\label{sec:grounding}

A grounding maps every symbol of the language to a real tensor or to an
operation on tensors \citep{badreddine2022ltn}. It starts by assigning to each
domain a set of tensors. The three subsections follow the three kinds of
expression: terms, predicates and formulas.

\subsection{Grounding of terms into real tensors}
\label{sec:grounding-terms}

A constant is grounded as one tensor. The tensor can be fixed, when the object
has known features, or it can be a learnable embedding
\citep{badreddine2022ltn}. A function symbol is grounded as a map between the
tensors of its domains. A variable is grounded as a sequence of individuals of
its domain, and a term with free variables is grounded as a tensor with one
axis for each variable. A term $g(x,y)$ then returns one result for each
combination of an individual of $x$ and one of $y$ \citep{badreddine2022ltn}.

In training, the sequence of a variable is a minibatch. A variable over
units is grounded as a batch of feature vectors, and the size of the batch
is the number of individuals over which a quantifier will range. The batch
size is therefore a design choice and not a detail of the implementation.

\subsection{Grounding of predicates into neural networks with output in [0,1]}
\label{sec:grounding-predicates}

A predicate is grounded as a function from the tensors of its input domains to
the interval $[0,1]$ \citep{badreddine2022ltn}. The value is a degree of
truth. This answers the question left open in \cref{ch:attacks-ddos}.
$\mathrm{HighRate}(x)$ no longer has to be true or false. For illustration, it
may return $0.9$ for a unit whose window has a very high rate and $0.2$ for a
modest one, and every axiom that contains it inherits the degree.

In practice a predicate is a neural network or another tensor operation whose
output lies in $[0,1]$. Badreddine et al. put a sigmoid on top of a multilayer
perceptron for a binary predicate. For a classifier with mutually exclusive
classes they use a softmax, so that the degrees of the classes sum to one
\citep{badreddine2022ltn}. Both forms are open to the thesis, because the
classes of the KDD families are exclusive for a record.

The thesis has two groups of predicates, as \cref{tab:rl-symbols} shows. The
class predicates are learned. The condition predicates are computed from the
features, either as fixed functions or as small learned ones. In our design the
network that grounds the class predicates has the architecture of the
multilayer perceptron baseline. Only the training objective differs, which is
the matched capacity requirement of \cref{sec:bias-variance-generalization}.

\subsection{Grounding of formulas into satisfaction degrees}
\label{sec:grounding-formulas}

Connectives and quantifiers are grounded with operators of fuzzy logic. A
connective is applied element by element to the tensors of its arguments. A
quantifier reduces the axis of the variable it binds with an aggregation
operator \citep{badreddine2022ltn}. Every formula therefore grounds to a
tensor with one axis for each free variable that remains, and a formula
without free variables grounds to a single number in $[0,1]$. This number is
its satisfaction degree. The operators are the subject of
\cref{sec:fuzzy-semantics}.

Take the SYN flood axiom and a batch of $B$ units. The three predicates
return three vectors of $B$ degrees. The conjunction and the implication
combine them position by position and give one vector of $B$ degrees. The
universal quantifier aggregates this vector into one number, the degree to
which the axiom holds on the batch. Learning, described in
\cref{sec:learning-satisfiability}, pushes this number up.

The same number, computed for one axiom on test data, is the raw material of
the axiom-level satisfaction of contribution C5 in \cref{sec:contributions}.
Its definition belongs to the benchmark, see
\cref{sec:detection-quality-dimension}.
% FIXME(open): \cref{sec:detection-quality-dimension} (Chapter 8, Sec. 8.2) renders as "??" until that chapter exists; C5 is defined formally there.

Variables of different domains can be quantified together. By default every
individual of one variable is paired with every individual of the others.
Diagonal quantification pairs instead the $i$-th individual of each variable
\citep{badreddine2022ltn}. This is the form that supervised axioms need, to
match an example with its own label. The observation is ours.


% ================================================================
% 5.3 Fuzzy semantics
% ================================================================
\section{Fuzzy semantics}
\label{sec:fuzzy-semantics}

The grounding of \cref{sec:grounding} leaves open which operators interpret
the connectives and the quantifiers. The choice matters for training.
\citet{vankrieken2021analyzing} show that many fuzzy operators, including
well-known ones, are unsuitable for gradient-based learning. Badreddine et al.
take up their results and recommend one configuration
\citep{badreddine2022ltn}. This section presents the operators, then the
gradient problems that decide between them. The configuration that the thesis
adopts is a design choice, and it can change the learning dynamics.

\subsection{T-norms and t-conorms: G\"odel, {\L}ukasiewicz, Product}
\label{sec:t-norms}

A fuzzy conjunction maps two degrees to one degree. It agrees with the Boolean
conjunction when both degrees are $0$ or $1$, and it does not decrease when an
argument grows. A t-norm is a conjunction that also satisfies $T(x,1)=x$ and
is commutative and associative \citep{badreddine2022ltn}. A t-conorm is the
same for disjunction, with $S(x,0)=x$. Given a negation $N$, the dual t-conorm
of $T$ is $S(x,y)=N(T(N(x),N(y)))$. A t-norm, its dual t-conorm and the
standard negation $N_S(a)=1-a$ form a symmetric configuration. There are three
common ones, shown in \cref{tab:t-norms}.

\begin{table}[htbp]
\centering
\footnotesize
\caption{The three common t-norms and their dual t-conorms
\citep{badreddine2022ltn}.}
\label{tab:t-norms}
\begin{tabular}{@{}lll@{}}
\toprule
Configuration & Conjunction $a\wedge b$ & Disjunction $a\vee b$ \\
\midrule
G\"odel & $\min(a,b)$ & $\max(a,b)$ \\
Product & $ab$ & $a+b-ab$ \\
{\L}ukasiewicz & $\max(a+b-1,0)$ & $\min(a+b,1)$ \\
\bottomrule
\end{tabular}
\end{table}

The three differ in the laws they keep. Only the pair of G\"odel operators is
distributive \citep{badreddine2022ltn}. For degrees in $\{0,1\}$ all three
return the Boolean value, so they differ only on graded degrees. This remark
is ours, and it matters because some conditions of our axioms may be crisp.

\subsection{Fuzzy implications: Kleene-Dienes, Goguen, {\L}ukasiewicz, Reichenbach}
\label{sec:fuzzy-implications}

Two classes of implication are built from the operators above
\citep{badreddine2022ltn}. An S-implication is defined as $\neg a\vee b$. An
R-implication is the largest $z$ with $a\wedge z\leq b$, which generalises
modus ponens. \Cref{tab:implications} lists the common ones, with the class and
the t-norm or t-conorm behind each.

\begin{table}[htbp]
\centering
\footnotesize
\caption{Common fuzzy implications \citep{badreddine2022ltn}.}
\label{tab:implications}
\begin{tabular}{@{}llll@{}}
\toprule
Name & $I(a,b)$ & Class & Based on \\
\midrule
Kleene-Dienes & $\max(1-a,b)$ & S & $S_M$ \\
G\"odel & $1$ if $a\leq b$, else $b$ & R & $T_M$ \\
Reichenbach & $1-a+ab$ & S & $S_P$ \\
Goguen & $1$ if $a\leq b$, else $b/a$ & R & $T_P$ \\
{\L}ukasiewicz & $\min(1-a+b,1)$ & S and R & $S_L$, $T_L$ \\
\bottomrule
\end{tabular}
\end{table}

Every axiom of \cref{ch:attacks-ddos} is an implication, so this choice weighs
more than the others. With Reichenbach the derivative of the implication with
respect to the consequent is the degree of the antecedent. A unit on which
the condition is weakly true pushes the class predicate weakly, and an axiom
acts mostly where its condition holds. We derive this from the formula. {\L}ukasiewicz and Reichenbach coincide when the
antecedent is crisp. Both return $b$ for $a=1$ and $1$ for $a=0$. They differ
only when the condition is graded, which is the case that the thesis adds.

\subsection{Aggregators for quantifiers: mean, min, generalised p-mean, pMeanError}
\label{sec:aggregators}

An aggregation operator maps any number of degrees to one degree. It is
monotone, returns $x$ for the single input $x$, and returns $0$ or $1$ for
inputs that are all $0$ or all $1$ \citep{badreddine2022ltn}. A t-norm applied
repeatedly gives a candidate for the universal quantifier, for example the
minimum or the product of all degrees. The product of many degrees below $1$
tends to $0$, which in our view makes it a poor choice for a batch of hundreds
of units.

Badreddine et al. propose two generalised means instead. For $n$ degrees
$a_1,\dots,a_n$ in $[0,1]$ and an exponent $p\geq 1$,
\begin{align}
\exists:\quad A_{pM}(a_1,\dots,a_n) &=
  \Bigl(\tfrac{1}{n}\textstyle\sum_{i=1}^{n} a_i^{p}\Bigr)^{1/p}, \\
\forall:\quad A_{pME}(a_1,\dots,a_n) &=
  1-\Bigl(\tfrac{1}{n}\textstyle\sum_{i=1}^{n}(1-a_i)^{p}\Bigr)^{1/p}.
\end{align}
$A_{pME}$ measures the deviation of each degree from $1$, and the two are
related by $\exists\equiv\neg\forall\neg$ \citep{badreddine2022ltn}. For $p=1$
both are the arithmetic mean. As $p$ grows, $A_{pM}$ approaches the maximum
and $A_{pME}$ the minimum. For $p=2$, $A_{pME}$ equals
$1-\mathrm{RMSE}(\mathbf{a},\mathbf{1})$, where $\mathbf{a}$ is the vector of
degrees. The exponent is a hyperparameter, and it can differ from formula to
formula \citep{badreddine2022ltn}. A formula $\forall x\,P(x)$ with a low $p$
says that $P$ holds for some $x$, and a higher $p$ says that it holds for most.

\subsection{Vanishing and exploding gradients, and the choice of p}
\label{sec:gradient-problems}

Three gradient problems arise in fuzzy operators \citep{badreddine2022ltn}. An
operator is single-passing if its derivative is non-zero for one argument
only. The gradient vanishes if it is null on part of the domain, and the
inputs there are not updated. The gradient explodes if large errors
accumulate and the updates become unstable. \Cref{tab:gradient-problems}
summarises them for the operators of the previous subsections.

\begin{table}[htbp]
\centering
\footnotesize
\caption{Gradient problems of some operators, from
\citet{badreddine2022ltn}. A parenthesis marks a problem that appears only on
an edge case.}
\label{tab:gradient-problems}
\begin{tabular}{@{}lccc@{}}
\toprule
Operator & Single-passing & Vanishing & Exploding \\
\midrule
$T_M$, $S_M$, $I_{KD}$ & yes & & \\
$I_G$ & yes & yes & \\
$T_P$, $S_P$, $I_R$ & & (yes) & \\
$T_L$, $S_L$, $I_{\L uk}$ & & yes & \\
min and max aggregators & yes & & \\
product and {\L}ukasiewicz aggregators & & yes & \\
$A_{pM}$, $A_{pME}$ & & & (yes) \\
\bottomrule
\end{tabular}
\end{table}

The table explains the recommended configuration. The product operators with
the Reichenbach implication fail only on edge cases. The edge cases are
$a=b=0$ for $T_P$, $a=b=1$ for $S_P$, and $a=0$, $b=1$ for $I_R$. $A_{pM}$
explodes when $\sum_i a_i^{p}$ tends to $0$, and $A_{pME}$ when
$\sum_i(1-a_i)^{p}$ tends to $0$ \citep{badreddine2022ltn}. The authors remove
these cases with two projections, where $\epsilon$ is an arbitrarily small
positive number,
\begin{align}
\pi_0(a) &= (1-\epsilon)\,a+\epsilon, \\
\pi_1(a) &= (1-\epsilon)\,a.
\end{align}
The first maps $[0,1]$ into $]0,1]$ and the second into $[0,1[$. The stable
operators apply them before the original ones, so that $T_P'(a,b)=
T_P(\pi_0(a),\pi_0(b))$ and $S_P'(a,b)=S_P(\pi_1(a),\pi_1(b))$. The
implication becomes $I_R'(a,b)=I_R(\pi_0(a),\pi_1(b))$, and the aggregators
are treated alike. Negation is unchanged. The result is the stable product
configuration \citep{badreddine2022ltn}.

The price is a small distortion. $T_P'(a,1)=(1-\epsilon)a+\epsilon\neq a$ for
$a<1$, so the stable conjunction is no longer a t-norm
\citep{badreddine2022ltn}. The authors report that the stable operators
improve numerical stability in their experiments. We take the product
configuration, with the Reichenbach implication and $A_{pME}$ for every
universal quantifier. Of the stable operators we keep the stable aggregator
only: the conjunction of the conditions and the Reichenbach implication are
computed in their plain form. $\epsilon$ is kept small and fixed, and it is not
varied in the experiments. The sensitivity experiment E4 varies the exponent
$p$, the crispness of the conditions and the scale of their widths
(\cref{sec:kb-format}).

The exponent $p$ decides where the gradient goes. The derivative of $A_{pME}$
with respect to $a_i$ contains the factor $(1-a_i)^{p-1}$
\citep{badreddine2022ltn}. For $p>1$ the units that satisfy the axiom least
receive the largest share of the gradient. A higher $p$ therefore concentrates
learning on the violations. The paper uses $p=2$ for the universal aggregator
in the binary, the multi-class and the digit addition examples. Its evidence on
scheduling $p$ concerns the existential quantifier in digit addition. There
$p=2$ from the first epoch gave poor results, and a schedule from $1$ to $6$
was the most promising approach \citep{badreddine2022ltn}. This evidence does
not transfer to the universal case that the thesis uses. We start from $p=2$
as in the paper and treat $p$ as a hyperparameter. A high $p$ has a cost in
our setting. An axiom such as the SYN flood one is a heuristic with
exceptions, and a single flash crowd could dominate its degree. This
argument is ours.
% FIXME(open): the value of p for each axiom is not fixed in Sec. 5.7; fix it in the methodology chapter (Ch. 8) with the other parameters of the logic.

% ================================================================
% 5.4 Learning as maximisation of satisfiability
% ================================================================
\section{Learning as maximisation of satisfiability}
\label{sec:learning-satisfiability}

The previous section gave every formula a degree of truth. This section turns
the degrees into a training objective.

\subsection{The knowledge base as a set of axioms}
\label{sec:knowledge-base}

A theory of Real Logic has three parts \citep{badreddine2022ltn}. The first is
a set of closed formulas, the knowledge base $\mathcal{K}$. The second is a
parametric grounding $\mathcal{G}(\cdot\mid\boldsymbol{\theta})$ of all the
symbols. The third is the hypothesis space of the parameters. The grounding
may be fixed for some symbols and learned for others. In the thesis the
condition predicates are mostly fixed and the class predicates are learned.

A formula aggregating operator $\mathrm{SatAgg}$ combines the degrees of all the
formulas of $\mathcal{K}$ into one number in $[0,1]$. This number is the
satisfiability of the theory \citep{badreddine2022ltn}. Learning searches for
the parameters that maximise it,
\[
\boldsymbol{\theta}^{*}=\operatorname*{argmax}_{\boldsymbol{\theta}\in\Theta}
\ \mathrm{SatAgg}_{\phi\in\mathcal{K}}\ \mathcal{G}_{\boldsymbol{\theta}}(\phi).
\]
A regularisation term $-\lambda\mathrm{R}(\boldsymbol{\theta})$ can be added,
as in ordinary machine learning \citep{badreddine2022ltn}. If the learned
symbols are predicates, the problem is a classification task. The thesis uses
only this case. The knowledge base has two groups of axioms. The supervised
axioms state the labels. The domain axioms are those of
\cref{ch:attacks-ddos}.

\subsection{Supervised axioms vs domain axioms; weighting the two groups}
\label{sec:supervised-domain-axioms}

In the multi-class example of Badreddine et al. the labels enter as axioms. A
variable is grounded with the examples of each class, and an axiom states that
the predicate of that class holds for them. Because the classes are mutually
exclusive, the output layer is a softmax, and exclusivity axioms are not
needed. The paper also contrasts one binary predicate $P(x,l)$ over a shared
network with one unary predicate per class. The unary predicates give disjoint
networks, and the binary predicate shares weights
\citep{badreddine2022ltn}. The thesis takes the shared form, one network whose
outputs are the class predicates, as in \cref{sec:grounding-predicates}.

The domain axioms are added to the same knowledge base. Training data and
constraints are then both formulas that must be satisfied
\citep{badreddine2022ltn}. The aggregator treats all formulas alike, so the
number of axioms in each group acts as an implicit weight. A weight can be
introduced explicitly. The authors note that it can be added to a formula
through a propositional predicate $p_w\wedge\phi$, whose value could then be
learned \citep{badreddine2022ltn}. The alternative is to aggregate each group
separately and combine the two numbers. The framework does not choose among
these. The weighting of the two groups is a design decision of the thesis.
% FIXME(open): the weighting of supervised formulas and security axioms is not fixed in Sec. 5.7; fix it in the methodology chapter (Ch. 8).

\subsection{Optimisation loop, batching and the role of the guarded quantifier}
\label{sec:optimisation-loop}

In practice the optimiser minimises
$L=1-\mathrm{SatAgg}_{\phi\in\mathcal{K}}\,\mathcal{G}_{\boldsymbol{\theta},
x\leftarrow B}(\phi)$, where $B$ is a minibatch drawn from the data
\citep{badreddine2022ltn}. The authors work out a binary classification with
two axioms, one for the positive and one for the negative examples. With the
stable product configuration and $p=2$ the loss reduces to a combination of
the root mean square errors of the two groups \citep{badreddine2022ltn}.
In our reading, supervised training with a logical loss is therefore ordinary
training with another error function. The objective depends on three choices, the operators,
their parameters such as $p$, and the formula aggregator. They are part of the
model definition, and the benchmark must record them.

The minibatch is not neutral. It sets the number of individuals that each
quantifier aggregates. The classes of the thesis are imbalanced, and a rare
class such as R2L or U2R can be missing from a batch. The variable of that
class is then empty. The paper handles empty masks by convention, with the
universal quantifier returning $1$ and the existential returning $0$
\citep{badreddine2022ltn}. The axiom is then trivially satisfied and, as we
infer, carries no gradient. The batch composition therefore decides whether a
rare class is trained at each step.

A guarded quantifier restricts the individuals before the aggregation, with a
Boolean mask \citep{badreddine2022ltn}. $\forall x:m(x)\,\phi(x)$ is not the
same as $\forall x\,(m(x)\rightarrow\phi(x))$. With the mean and Reichenbach,
a mask of $(0,1,1)$ over degrees $(0.2,0.7,0.8)$ gives $0.75$ for the guarded
form and $0.833$ for the implication, and the guarded form has no gradient for
the masked-out individuals \citep{badreddine2022ltn}. The axioms of the
thesis are written as implications, so they do not use guards. Guards appear
in a different place. A variable grounded with the examples of one class is
the same object as a quantifier guarded by the label, and it is how the
supervised axioms of the previous subsection are built.

% ================================================================
% 5.5 What LTN can and cannot do
% ================================================================
\section{What LTN can and cannot do}
\label{sec:ltn-capabilities}

The sections so far describe how the framework works. This section says what it
gives the thesis and where it stops.

\subsection{Querying: satisfaction of a formula as a first-class output}
\label{sec:querying}

A trained theory can be queried \citep{badreddine2022ltn}. A truth query
returns the degree of a formula. A generalisation truth query evaluates a
formula with a free variable on unseen individuals and returns one degree for
each. The paper defines two more query types, for terms. The generalisation truth
query is the one that the benchmark needs. The degree of an axiom on a test set is computed
without retraining, and it can be computed for any axiom, also for one that was
not in the training knowledge base. The per-unit degrees come before the
aggregation, so they also show which units violate the axiom and not only how
many. This use is ours, and it supports the explanation question of RQ4.

A degree is not a guarantee, for two reasons. The first is the aggregator.
With $p=2$ the degree of a universal formula is $1-\mathrm{RMSE}$, and very
different situations give the same number. Ten percent of units that violate
an axiom completely and half of the units with a degree of about $0.55$ both
give about $0.68$. The example is ours. A single number hides the distribution,
and the benchmark must report the operators and also a rate of violations.
% FIXME(open): the rate of violations is the domain-rule violation rate of Ch. 8; \cref{sec:detection-quality-dimension} must say which threshold turns a degree into a violation.
The second reason is deeper. Badreddine et al. define logical consequence with
a threshold, and they show that querying after learning can accept a false
consequence. In their example every run of the optimiser reaches the same
optimum, no counterexample is found, and the consequence is wrongly assumed to
hold \citep{badreddine2022ltn}. Their remedy is a search for counterexamples,
which the thesis does not use. More generally, methods that put the constraints
in the loss cannot guarantee that the constraints are satisfied, although they
can generalise from less data or use unlabelled data \citep{giunchiglia2022deep}.
The thesis therefore measures satisfaction and does not assume it.

\subsection{Reasoning at training time vs reasoning at inference time}
\label{sec:inference-vs-training}

In a constraint-based model the logic acts in the loss. At each training step
the loss checks whether the scores satisfy the formulas. At inference only the
scores of the network are used for the prediction, and the formulas are not
evaluated \citep{marra2024statistical}. This is the contrast with proof-based
systems, where a logic program is part of the model. Badreddine et al. say the
same for the cost. In LTN, inference involves only a forward pass through the
neural component \citep{badreddine2022ltn}.

For the thesis the consequence is the following. The class predicates are the
outputs of one network with the architecture of the baseline. The axioms shape
its parameters during training. At inference the network is called as a plain
multilayer perceptron, with the same number of parameters and the same forward
pass as the baseline. The condition predicates and the axioms are not part of
the prediction. This is our reading of the sources for our design, and the
latency experiment E12, defined in \cref{sec:symbolic-placement}, is meant to test it.
It has a limit. The rules are not enforced at inference. A unit that violates
an axiom is still classified by the network, which is why the satisfaction on
test data is a result and not a property. Keeping a symbolic component in the
loop at inference would be a different design with a different cost, for
example a rule check after the prediction. The thesis does not make it one of
the compared arms, and measures it as an optional post-check in the second part
of E12.

\subsection{Scalability limits with many axioms and large batches}
\label{sec:scalability-limits}

The authors report that training and test times of LTN are of the same order
of magnitude as those of a CNN baseline on digit addition. The measured steps
take $5.36$ against $2.72$ ms in training and $3.44$ against $1.45$ ms in
test, on one GPU with batches of $32$ \citep{badreddine2022ltn}. By our
calculation LTN takes about twice as long in training and about 2.4 times as
long in test. The task is image
classification, and nothing is shown for a tabular multilayer perceptron, so
the figure is a warning and not a prediction. The cost of the thesis models is
measured in the experiments.

Two limits follow from the grounding. A formula with several free variables
grounds to a tensor with one axis for each, and by default every individual of
one variable is paired with every individual of the others, so the size is the
product of the batch sizes \citep{badreddine2022ltn}. Diagonal quantification
avoids it when examples match one to one. The axioms of \cref{ch:attacks-ddos}
have a single free variable, and sources are described by features computed
beforehand, so the thesis stays on the cheap side. The second limit concerns the
number of axioms. If every class predicate wraps the whole network, each use of
a predicate in an axiom runs the forward pass again. The cost then grows with
the number of occurrences. Computing the outputs once per batch and letting the
predicates read them avoids this. The point is ours, and it belongs to the
implementation.

% ================================================================
% 5.6 Alternatives and why LTN is chosen here
% ================================================================
\section{Alternatives and why LTN is chosen here}
\label{sec:ltn-alternatives}

\subsection{LTN vs Semantic Loss vs DeepProbLog vs NeurASP: criteria and trade-offs}
\label{sec:ltn-vs-alternatives}

We compare the four systems on three criteria. The first is the semantics of
the logic. The second is the role of the logic, as a term of the loss or as a
program that the network feeds. The third is what happens at inference.
\Cref{tab:ltn-alternatives} gives the result.

\begin{table}[htbp]
\centering
\footnotesize
\setlength{\tabcolsep}{3pt}
\caption{LTN and the three alternatives on the criteria of this section.
Sources: \citet{badreddine2022ltn,marra2024statistical,xu2018semantic,%
manhaeve2018deepproblog,yang2020neurasp}.}
\label{tab:ltn-alternatives}
\begin{tabular}{@{}p{2.1cm}p{3.3cm}p{3.5cm}p{3.1cm}@{}}
\toprule
System & Semantics & Role of the logic & At inference \\
\midrule
LTN & Fuzzy, first-order, many-valued & Term of the loss & Network only \\
Semantic Loss & Probabilistic, propositional constraints & Term of the loss & Network only \\
DeepProbLog & Probabilistic logic program & Network feeds a program & Program is part of the model \\
NeurASP & Answer set program with probabilities & Network feeds a program & Program is part of the model \\
\bottomrule
\end{tabular}
\end{table}

The semantic loss compiles propositional constraints into arithmetic circuits
for weighted model counting, and the loss measures how close the network is to
satisfying them \citep{badreddine2022ltn}. DeepProbLog
reads the output of a network as the probability of an atom. A query is
answered by probabilistic inference based on model counting, with a closed
world assumption and a language without function symbols. The clauses are
compiled into a circuit, and the compilation is costly. In LTN the truth value
of a formula is computed from the values of its atoms, which the authors judge
more efficient than model counting and compatible with full first-order logic
\citep{badreddine2022ltn}. The two kinds of inference are not comparable
\citep{badreddine2022ltn}. NeurASP reads the outputs of a network as
probability distributions over the facts of an answer set program
\citep{yang2020neurasp}.

We choose LTN for four reasons, and they are our judgment. First, the
architecture stays that of the baseline. The comparison of the thesis isolates
the effect of knowledge in the loss, and a program that the network feeds would
make the model something else than a multilayer perceptron. Second, the
conditions of the axioms are graded, such as a high rate, and a degree of truth
expresses that. A probabilistic reading would ask for the probability of a
condition that has no frequency behind it. Third, the cost at inference stays
that of the baseline, in line with the use case of the thesis. Fourth, the
framework has a TensorFlow 2 implementation \citep{badreddine2022ltn}, the
framework chosen for the work.

The choice has limits that we state. The sources read here do not compare the
four systems on network traffic, and the thesis does not run the other three.
Neural predicates could also adapt to rules over predicates, as the previous
chapter notes, and a program that derives the attack class from scored
conditions would put the logic in the decision. That design tests another
question. The choice of LTN is therefore a choice of what to measure, and not a
proof that it is the best system. The comparison of tooling maturity is left
open.
% FIXME(open): the tooling maturity of DeepProbLog, DeepStochLog, NeurASP, DL2 and semantic loss is not compared; do it before Ch. 9 or declare it a limitation in Ch. 11.

\subsection{Place in the taxonomy of Kautz}
\label{sec:kautz-place}

\Cref{sec:kautz-taxonomy} left the type of the LTN models to this chapter.
Kautz names logic tensor networks as an example of the fifth type
\citep{kautz2022third}, and \citet{wan2024cognitive} describe that type as soft
constraints or regularizers. The survey of \citet{marra2024statistical} places
the framework among the systems that turn formulas into a regularisation loss.
We keep the assignment to the fifth type for the framework. The thesis uses
the framework at the level of the loss only. The architecture of the network
is not built from the rules. The description of the fourth type in
\cref{tab:kautz-types}, a training regime based on symbolic rules, fits this
use at least as well as the description of the fifth, symbolic rules as
templates for structures. The assignment is therefore not confirmed without
reservation, and we call our models the fifth type by the authority of the
example, with this reservation.

% ================================================================
% 5.7 Formalising network security knowledge in Real Logic
% ================================================================
\section{Formalising network security knowledge in Real Logic}
\label{sec:formalising-security-knowledge}

This section writes the knowledge of \cref{ch:attacks-ddos} in the language
of the chapter. It settles what the framework leaves open, which are the
individuals, the shape of the condition predicates and the axiom set.

\subsection{Predicate design over flow features}
\label{sec:predicate-design}

The classifier labels one unit for each decision. In NSL-KDD it is a
connection record, in the flow datasets it is a flow. The axioms of
\cref{tab:candidate-axioms} speak of windows, flows and sources. We attach
to every unit the statistics of its window and of its source. For the window
of a connection record, the traffic features of
\cref{sec:time-window-connection-aggregates} already do this, but they group
connections by destination host or service. NSL-KDD records carry no source
address, so a statistic over the connections of one source cannot be
computed there. In the flow datasets the window and source statistics are
not given and must be computed from the flows. All predicates then take one variable $x$ over
units, and $f$, $w$ and $s$ name only the scope of a statistic. The slow DoS
axiom, for example, reads a flow statistic of $x$, and the SYN flood axiom
reads the statistics of the window of $x$. We take this design because the
classes are predicates of the labelled unit, and because it avoids a
membership relation between flows, windows and sources inside the logic.

The design has a cost. Many units share the statistics of one source, so an
axiom on sources is evaluated once for each unit and not once for each source.
A source with many units weighs more in the satisfaction degree, a
consequence we draw from the design.

A condition predicate reads one or two of these statistics and returns a
degree. The class predicates are the outputs of the network, as in
\cref{sec:grounding-predicates}. The mapping of each condition to the
features of each dataset belongs to the implementation.
% FIXME(open): the mapping of every condition predicate to the features of NSL-KDD, CIC-IDS2017 and CIC-DDoS2019 is given in Ch. 9, \cref{ch:implementation}.

\subsection{Discretisation and fuzzification of continuous features: thresholds vs sigmoids}
\label{sec:fuzzification}

A condition such as a high rate needs a threshold. Two forms are possible. A
crisp predicate returns $1$ above the threshold $\tau$ and $0$ below it. A
graded predicate returns $\sigma\bigl((v-\tau)/s\bigr)$ for a feature $v$,
where $\sigma$ is the logistic function and $s>0$ sets the width of the
transition. The crisp form is the limit of the graded one for $s\to 0$. The
choice has two effects. A crisp condition gives a unit just under the
threshold no influence, and one just over it full influence. A graded
condition gives each unit a partial influence, and with the Reichenbach
implication of \cref{sec:fuzzy-implications} this is a gradient weight on the
consequent. We prefer the graded form for this reason, and we keep the crisp
one as a baseline.

The parameters $\tau$ and $s$ are hyperparameters of the logic. They must be
chosen on the training and validation splits and never on the test set. They
should not be learned together with the network. An implication is true when
its antecedent is false, so a learnable threshold lets the optimiser satisfy
an axiom by moving the threshold until the condition never holds. The
satisfaction degree rises and the axiom no longer says anything. This risk is
ours, and it is one more reason to fix the thresholds.

\subsection{The axiom set inherited from the internship, restated formally}
\label{sec:axiom-set-internship}

The thesis builds on the internship of its author, whose report is not
published. The report defines four domain axioms on NSL-KDD, guided by an
analysis of the errors of the model on the training data. Written in Real
Logic, with the conditions of the report, they are
\begin{align}
&\forall x\,\bigl(\mathrm{DiffSrv}(x)\vee\mathrm{DstHostDiffSrv}(x)\rightarrow
  \mathrm{Probe}(x)\bigr),\\
&\forall x\,\bigl(\mathrm{GuestLogin}(x)\vee\mathrm{FailedLogins}(x)\vee
  \mathrm{Hot}(x)\rightarrow\mathrm{R2L}(x)\bigr),\\
&\forall x\,\bigl(\mathrm{RootShell}(x)\vee\mathrm{SuAttempted}(x)\vee
  \mathrm{FileCreations}(x)\rightarrow\mathrm{U2R}(x)\bigr),\\
&\forall x\,\bigl(\mathrm{LoggedIn}(x)\wedge\mathrm{NoSynErrors}(x)\wedge
  \mathrm{NoHot}(x)\rightarrow\mathrm{Normal}(x)\bigr).
\end{align}
The first axiom reads the rates of connections to different services, and the
last one protects legitimate traffic. The report uses crisp conditions on
standardised features, mostly a threshold of $0.5$, and the binary features
as they are. We restate them here with graded predicates, and the
thresholds are open to the choice of the previous subsection.

Two points limit their use. The conditions of the R2L and U2R axioms are
content features of NSL-KDD, such as a guest login or a root shell. Flow
records do not carry them, as \cref{sec:r2l-u2r-hard} notes. These axioms can
be evaluated on NSL-KDD only, and a detector deployed on flow records could not
use them. This is our judgment, and the benchmark keeps them, one axiom for each term of
the disjunctions, in a variant of the NSL-KDD knowledge base with weak knowledge
(\cref{sec:kb-format}). Second, the axioms were written after looking at the errors of
one model on the training data. They are tuned to that model and dataset, and
they are not independent knowledge like those of \cref{ch:attacks-ddos}. The
results of the report are not used in this thesis.

\subsection{Extending it: DDoS-specific axioms, mutual exclusivity, hierarchy axioms}
\label{sec:ddos-axioms}

The seven axioms of \cref{tab:candidate-axioms} are the DDoS-specific group,
written with the form of \cref{sec:predicate-design}. They cover the class that
the internship axioms leave without a rule. Most of them need a window or a
source, and \cref{sec:kb-elicitation} says which ones the knowledge bases of the
experiments can hold.

Two structural groups can be added. Mutual exclusivity states that a unit has
one class, for example $\forall x\,\neg(\mathrm{Normal}(x)\wedge\mathrm{U2R}(x))$.
\citet{onchis2022neurosymbolic} use axioms of this form in a classifier built
with logic tensor networks and applied to KDD99 and CIC-IDS2017. With a softmax
output they are redundant, as the exclusivity is then built into the grounding
\citep{badreddine2022ltn}. They matter only if the classes use independent
sigmoids. Hierarchy axioms relate a subclass to its class, for example
$\forall x\,(\mathrm{Reflection}(x)\rightarrow\mathrm{DoS}(x))$ and
$\forall x\,(\mathrm{DoS}(x)\rightarrow\neg\mathrm{Normal}(x))$. They need a
dataset with the finer labels. CIC-DDoS2019 has reflection attacks over several
protocols \citep{cic2019ddos}. Whether the structural axioms add anything to
a softmax is an empirical question of the benchmark.

% ================================================================
% 5.8 Summary
% ================================================================
\section{Summary}
\label{sec:ch5-summary}

Real Logic gives the predicates of \cref{ch:attacks-ddos} a degree of truth, and
a fuzzy semantics turns a formula into a number that a network can maximise.
The thesis takes the product configuration with a stable aggregator and
trains a network with the architecture of the baseline on supervised and domain axioms. At inference
the network is a plain multilayer perceptron, and the satisfaction of an axiom
is measured and not assumed. The candidate axioms are the seven of the attack
taxonomy and the four of the internship, over one domain of classified units with
graded conditions, and \cref{ch:usecases-design} draws the knowledge bases of the
experiments from them and from the training data. The next chapter reviews the work on neurosymbolic detection
that came before.